\section{Results: a recombined PETase library, assayed}
\label{sec:res:experimental}

\plan{\textbf{Claim.} A fragment-recombination library built by exactly the
procedure of Section~\ref{sec:res:fragments} was ordered, assembled and assayed:
3{,}135 designs, each 18--76 mutations from its parent, measured on a single
run of the MoMI esterase assay. Two things follow. The platform produces
physically assemblable, expressible, measurable libraries at combinatorial
scale --- and zero-shot structure prediction, the signal the policy is trained
against, is \emph{lightly} predictive of measured activity: $\rho = 0.58$ at
best, worth a $2.8\times$ enrichment, and beaten by mutation count alone.

\textbf{Status: \statusdone{}.} Every number below is on disk and is regenerated
by \texttt{paper/figures/petase\_zeroshot.py} from
\texttt{figures/data/petase\_activity\_zeroshot.csv}; the figure script prints
the prose's numbers so the two cannot drift. What this section does \emph{not}
contain is a closed-loop activity campaign --- see
``What this does not show'' at the end, and the plan in
\texttt{archive/results\_petase.tex}.}

\subsection*{The library is the method, physically}

The assayed library is a five-fragment Golden Gate recombination library over a
209-residue de novo PETase scaffold with the Ser66/His107/Asp122 triad.
Boundaries fall at residues giving fragments of 42, 44, 42, 37 and 44 residues.
The curated fragment dictionary carries 25, 26, 26, 26 and 25 variants at the
five slots: \textbf{128 protein parts spanning $1.1 \times 10^{7}$
recombinants}, and across library generations 1{,}920 DNA parts were ordered as
BsaI-flanked blocks and assembled one-pot from an acoustic-transfer map, one
well per recombinant. This is the arithmetic of Section~\ref{sec:res:fragments}
instantiated in DNA: the ratio of orderable parts to reachable variants is
roughly $10^{5}$.

What was recovered and assayed is 3{,}135 designs carrying a median of 29
mutations from the parent (interquartile range 25--38, full range 18--76).
\textbf{67\% register rate above background}, which is the first result: a
library recombined across parents at this mutational depth is mostly not dead.
That is the assumption Section~\ref{sec:res:fragments} names as load-bearing --
that fragments contribute roughly independently to foldability -- and it is the
only direct experimental test of it in this paper.

\subsection*{Zero-shot prediction is lightly predictive}

\begin{figure}[t]
  \centering
  \includegraphics[width=\textwidth]{figures/pdf/petase_zeroshot.pdf}
  \caption{\textbf{Structure-prediction confidence against measured activity
  for 3{,}135 recombined PETase designs.}
  (a) The strongest zero-shot signal is local and state-specific: per-atom
  pLDDT over the active site of the \emph{acyl-enzyme intermediate}. Selecting
  the top 5\% by this metric returns top-decile actives at $2.8\times$ the base
  rate. (b) Activity by design round, on the single screen run. (c) The
  control: distance from the design parent ranks activity better than any
  prediction does. The 20 open red points are a selected set re-assayed in a
  different run of the same assay and are never pooled into a correlation.}
  \label{fig:petase-zeroshot}
\end{figure}

Every design was folded by AF3 in two states --- the Michaelis complex and the
covalent acyl-enzyme intermediate --- and scored on 52 confidence and geometry
metrics. Table~\ref{tab:petase-zeroshot} gives the ranking against measured
activity.

\begin{table}[h]
  \centering
  \small
  \begin{tabular}{llrr}
    \toprule
    Predictor & What it measures & Spearman $\rho$ & $n$ \\
    \midrule
    \multicolumn{4}{l}{\emph{Zero-shot, AF3}} \\
    \quad acyl His107-NE2 pLDDT      & catalytic atom, acyl state & $+0.60$ & 2{,}548 \\
    \quad acyl active-site pLDDT     & active-site shell, acyl state & $+0.58$ & 2{,}548 \\
    \quad acyl PAE Ser66$\leftrightarrow$His107 & triad geometry & $-0.47$ & 2{,}548 \\
    \quad substrate interface PAE (min) & protein--ligand placement & $-0.43$ & 2{,}616 \\
    \quad substrate pTM               & global confidence & $+0.36$ & 2{,}616 \\
    \quad substrate ipTM              & global interface confidence & $+0.32$ & 2{,}616 \\
    \midrule
    \multicolumn{4}{l}{\emph{Zero-shot, sequence only}} \\
    \quad mutations from parent      & distance, no structure at all & $-0.65$ & 3{,}135 \\
    \midrule
    \multicolumn{4}{l}{\emph{Supervised, held out} (RidgeCV, 200 stratified splits)} \\
    \quad AF3 features               & & $+0.73 \pm 0.02$ & \\
    \quad Rosetta features           & & $+0.75 \pm 0.01$ & \\
    \quad AF3 $+$ Rosetta            & & $+0.80 \pm 0.01$ & \\
    \bottomrule
  \end{tabular}
  \caption{\textbf{Ranking power against measured activity.} Zero-shot rows are
  univariate Spearman on the screen run. Global confidence is roughly half as
  informative as a local, state-specific measurement at the catalytic atoms.}
  \label{tab:petase-zeroshot}
\end{table}

Three readings, in the order that matters.

\textbf{The useful signal is local and state-specific, not global.} pTM and ipTM
--- the metrics a design pipeline reaches for by default --- reach $\rho \approx
0.35$. Per-atom pLDDT at the catalytic histidine of the \emph{acyl-enzyme
intermediate} reaches $0.60$. The gain comes from scoring the chemically
relevant state at the chemically relevant atoms, which is an argument for
building the covalent intermediate into the oracle rather than for folding
harder.

\textbf{It is a real signal, and it is light.} Selecting the top 5\% of the
library by active-site pLDDT returns a top-decile-active design 28\% of the
time against a 10\% base rate: a $2.8\times$ enrichment, falling to $1.9\times$
if 25\% of the library is kept. That is worth having when the alternative is
random ordering, and it is nowhere near a filter one would trust to replace an
assay.

\textbf{And a cheaper predictor beats it.} Mutation count from the parent, which
requires no structure prediction at all, ranks activity at $\rho = -0.65$ ---
better than any AF3 metric in the table (Figure~\ref{fig:petase-zeroshot}c).
This is the comparison a zero-shot claim has to clear and it does not clear it.
What survives is that the AF3 signal is not merely that confound re-expressed:
restricted to the 2{,}539 designs at least 20 mutations from the parent, active-site
pLDDT still gives $\rho = +0.59$, and within single 5-mutation bands it ranges
from $+0.07$ to $+0.67$. Prediction adds something to distance; it does not
dominate it.

\subsection*{Searching the library moves activity}

Across four design rounds on the single screen run, median activity rises from
2.2 to 12.2 and the best recombinant from 20.5 to 47.8
(Figure~\ref{fig:petase-zeroshot}b). The progression is not monotone --- round 2
has the \emph{lowest} median of the four and the \emph{best} single design,
which is what a deliberately broadened library looks like.

A separate, selected set of 20 designs was re-assayed and scored 204--261,
roughly 100$\times$ the screen median. \textbf{Those numbers are from a different
run of the same assay} and the two runs share no normalizing control in the
record we have, so this paper does not quote a fold-improvement that crosses
that line. The 20 also sit 1--11 mutations from the parent, not out in the
recombined library, so they are evidence about a local optimization around a
good design --- not about what recombination reached.

\subsection*{What this does not show}

\begin{planlist}
  \item \textbf{No wild-type comparator.} The parent sequence was never assayed
        in this table, so no fold-over-parent number can be computed from it,
        and none is quoted.
  \item \textbf{Activity falls with distance.} The most active designs in the
        record are the closest to the parent. Recombination at 20--80 mutations
        buys reachable diversity and a 67\% register rate; on this scaffold and
        this assay it did not buy a better enzyme than careful local work did.
        Section~\ref{sec:res:frontier}'s claim is about the diversity/fidelity
        frontier, and that is the claim this section supports.
  \item \textbf{One scaffold, one chemistry, one assay.} The AME results carry
        the generality claim; this section carries the physical-realizability
        claim and the calibration of the oracle.
  \item \textbf{The closed loop is not here.} The surrogate-in-the-reward
        campaign is planned in \texttt{archive/results\_petase.tex} and has no
        data in this repository.
\end{planlist}
